\chapter{Testing \& Validation}
\label{ch:testing}

Rigorous testing was conducted to ensure the reliability, accuracy, and robust performance of the Mountain Climber Health \& GPS Tracker system under simulated field conditions. This chapter presents the methodologies and results of various hardware and software tests.

\section{LoRa Range Testing}
\label{sec:lora_range}

Given the critical nature of communication in remote areas, the SX1278 LoRa modules were tested extensively for range and signal integrity under different spreading factors (SF) and line-of-sight (LOS) conditions.

\begin{infobox}
Testing was conducted at 433MHz with a bandwidth of 125kHz and a coding rate of 4/5. Transmission power was set to 20 dBm.
\end{infobox}

\subsection{Test Methodology}
Two nodes were utilized: a transmitter periodically sending 50-byte packets and a receiver logging the Received Signal Strength Indicator (RSSI) and Packet Delivery Ratio (PDR). Tests were performed in an open field (LOS) and a moderately wooded suburban area (Non-LOS).

\subsection{Results}

\begin{table}[htbp]
    \centering
    \caption{LoRa Range Testing Results (Transmitter Power: 20 dBm, BW: 125kHz)}
    \label{tab:lora_results}
    \resizebox{\textwidth}{!}{
    \begin{tabular}{llccc}
        \toprule
        \textbf{Environment} & \textbf{Spreading Factor} & \textbf{Distance (km)} & \textbf{Avg RSSI (dBm)} & \textbf{Packet Loss (\%)} \\
        \midrule
        Line-of-Sight (LOS)  & SF7  & 1.5 & -95  & 0.0 \\
        Line-of-Sight (LOS)  & SF7  & 3.0 & -108 & 2.1 \\
        Line-of-Sight (LOS)  & SF10 & 5.0 & -115 & 1.5 \\
        Non-LOS (Wooded)     & SF7  & 0.5 & -102 & 4.2 \\
        Non-LOS (Wooded)     & SF7  & 1.2 & -118 & 22.5 \\
        Non-LOS (Wooded)     & SF10 & 1.5 & -112 & 3.0 \\
        Non-LOS (Wooded)     & SF12 & 2.5 & -120 & 8.5 \\
        \bottomrule
    \end{tabular}
    }
\end{table}

The results in Table \ref{tab:lora_results} indicate that while SF7 provides adequate LOS range, utilizing higher spreading factors like SF10 or SF12 is crucial for maintaining low packet loss in non-LOS mountainous environments, albeit at the cost of lower data rates.

\begin{figure}[htbp]
    \centering
    \includegraphics[width=0.8\textwidth]{lora_test_setup.png}
    \caption{LoRa Range Testing Setup in Field Conditions}
    \label{fig:lora_test_setup}
\end{figure}

\section{Sensor Accuracy Validation}
\label{sec:sensor_testing}

\subsection{MAX30102 Health Sensor}
The MAX30102 sensor's heart rate (BPM) and blood oxygen saturation (SpO2) readings were validated against a medical-grade finger pulse oximeter (ChoiceMMed MD300C2). 

\begin{table}[htbp]
    \centering
    \caption{MAX30102 Accuracy Comparison (Resting State, N=50 samples)}
    \label{tab:max_accuracy}
    \resizebox{\textwidth}{!}{
    \begin{tabular}{lccc}
        \toprule
        \textbf{Metric} & \textbf{Reference Device (Avg)} & \textbf{MAX30102 (Avg)} & \textbf{Mean Absolute Error} \\
        \midrule
        Heart Rate (BPM) & 72.4 & 73.1 & 1.8 BPM \\
        SpO2 (\%)        & 98.2 & 97.5 & 0.9 \% \\
        \bottomrule
    \end{tabular}
    }
\end{table}

\begin{warningbox}
While resting accuracy is within acceptable margins, vigorous wrist movement introduces motion artifacts, occasionally resulting in SpO2 reading dropouts.
\end{warningbox}

\subsection{NEO-6M GPS Module}
GPS accuracy was tested by placing the module at known surveyed coordinates. The module achieved a typical horizontal position error of 2.5 to 4 meters under clear skies (Time-To-First-Fix: $\sim$35 seconds from cold start).

\section{BLE Connection Reliability}
\label{sec:ble_testing}

The BLE link between the Wristband/Climber Main device and the Flutter mobile application was tested for connection stability and automatic reconnection.

\begin{itemize}
    \item \textbf{Connection Time:} Average 1.2 seconds to establish connection and discover GATT services.
    \item \textbf{Range:} Reliable connection maintained up to 8 meters indoors and 15 meters outdoors.
    \item \textbf{Reconnection Behavior:} When the connection drops due to distance, the Flutter app successfully reconnects automatically within 3 seconds of re-entering the range.
\end{itemize}

\section{System Integration Testing}
\label{sec:integration_testing}

An end-to-end integration test was conducted to verify the complete data pipeline from the climber to the base camp dashboard.

\textbf{Test Scenario: SOS Emergency Trigger}
\begin{enumerate}
    \item \textbf{Trigger:} Climber presses the physical SOS button on the Main Device.
    \item \textbf{Local Action:} Device buzzer activates, RGB LED turns red, and the Flutter app immediately displays a local emergency warning via BLE.
    \item \textbf{LoRa Transmission:} An SOS packet (Type 0x01) containing current GPS coordinates is broadcast via LoRa.
    \item \textbf{Relay:} The intermediate Repeater node successfully receives, validates the CRC, and re-broadcasts the packet.
    \item \textbf{Reception:} Base Camp unit receives the packet, outputs JSON over serial to the connected laptop.
    \item \textbf{Dashboard Alert:} The Flask dashboard parses the JSON, plays a loud audio alarm, flashes the screen red, and updates the climber's marker on the real-time map to indicate an emergency state.
\end{enumerate}

\textbf{Result:} The entire pipeline, from button press to dashboard alarm, consistently executed in under 1.5 seconds.
